\section{首轮产品队列服务于比较，不构成排他菜单}

现有十只产品被完整列入首轮研究，是因为它们曾是策略研究的起点，并不意味着它们构成长期唯一选择。开放研究宇宙允许直接股票、ETF、公募REIT、债券和现金工具在同一轮决策中竞争；但每一次组合评审都必须先冻结有限候选名单。一个未命名的“类别”不能获得估值、总回报或分散化信用，一个基金标签也不能替代底层经济暴露的研究。

\begin{exhibitbox}[首轮已尽调产品队列：用于后续跨资产比较，当前均为$q=0$]
\footnotesize
\begin{tabularx}{\exhibitboxwidth}{L{1.35cm}L{2.45cm}L{2.3cm}P}
\toprule
\textbf{代码} & \textbf{产品} & \textbf{已确认的研究角色} & \textbf{进入T3前仍须完成的事项} \\
\midrule
\texttt{511010} & 五年期国债ETF & 利率曲线暴露，不是受保障流动性。 & 同日价格/参考值、久期/DV01、价差、深度、容量、账户及组合压力。 \\
\texttt{510880} & 上证红利ETF & 红利权益暴露；采矿和金融权重会形成周期与金融因子集中。 & 基金当前成份、TTM实收现金、税后现金桥、发行人集团/因子/REIT交叉风险与容量。 \\
\texttt{512890} & 红利低波ETF & 低波名称不等于低组合风险；金融权重高且近年基金层分配有限。 & 基金当前成份、分配政策、现金质量、与其他权益候选的重叠、估值与容量。 \\
\texttt{508069} & 南京交通REIT & 单项目收费公路现金流。 & 车流、收费、期限、债务、当前估值、流动性和共同风险。 \\
\texttt{508011} & 物美消费REIT & 社区商业租金与租户现金流。 & 租户信用、租约、债务、资本开支、NAV日期与可执行容量。 \\
\texttt{508068} & 北京保障房REIT & 政策租金、出租和收缴。 & 扩募后可比估值、运营、债务、税费和组合层共同风险。 \\
\texttt{508077} & 有巢租赁住房REIT & 租赁住房现金流。 & 同簇住房REIT的相关性、经营KPI、估值、债务和流动性。 \\
\texttt{508055} & 上海租赁住房REIT & 租赁住房现金流。 & 地区/管理人/租赁因子重叠、运营、债务、估值和容量。 \\
\texttt{508080} & 亦庄产业园REIT & 产业园租户现金流；历史价格/NAV日期错配诊断为负。 & 租户集中、出租率、续约、债务、现金流、当前估值和流动性。 \\
\texttt{508096} & 中航京能国际能源REIT & 水光能源现金流。 & NAV口径勾稽、发电/电价、应收与保理、债务、估值和流动性。 \\
\bottomrule
\end{tabularx}
\end{exhibitbox}

这十项已具备产品级资料卡，不等于它们优于直接股票或其他合规工具。对直接股票，后续T2名单中的每一家公司也必须按\texttt{equity-research}标准完成独立深度报告；对ETF和REIT，则用结构适配的产品尽调替代公司模板，而不是降低证据标准。

\section{公募REIT：具有分配机制的研究候选，而非回报承诺}

基础设施REIT在符合分配条件时，规则要求每年将合并后年度可供分配金额的90\%以上以现金形式分配；可供分配金额须在净利润基础上调整，并综合考虑持续发展、偿债能力和经营现金流。该规则说明产品机制，不规定买入成本上的分配率，也不保证二级市场总回报。\sourcenote{公开募集基础设施证券投资基金指引（试行）（R01）。}

\begin{exhibitbox}[已披露现金事实不等于可配置结论]
\footnotesize
\begin{tabularx}{\exhibitboxwidth}{L{1.4cm}L{2.6cm}L{2.2cm}P}
\toprule
\textbf{代码} & \textbf{候选} & \textbf{已披露现金事实} & \textbf{为何仍不能以此批准配置} \\
\midrule
\texttt{508069} & 华夏南京交通高速公路REIT & 2025年化现金流分派率9.19\%；2026年首期每10份CNY1.4980。同期年报披露内部收益率1.44\%。 & 当前价格、TTM分配、收费量、剩余期限、债务、估值和流动性缺失。 \\
\texttt{508011} & 嘉实物美消费REIT & 2025-12-31基准日可供分配CNY35.8781m，拟分配CNY35.2400m。 & 当前市价/NAV、消费物业运营、租户、债务与成交容量缺失。 \\
\texttt{508068} & 华夏北京保障房REIT & 2025年合并可供分配CNY74.8239m；本次每10份CNY0.6000。 & 当前市价/NAV、出租率、租金、资本开支和集中度缺失。 \\
\texttt{508077} & 华夏华润有巢租赁住房REIT & 2025年可供分配CNY50.5599m；本次每10份CNY0.1462。 & 与其他租赁住房候选的行业重叠、当前估值和流动性缺失。 \\
\texttt{508055} & 汇添富上海地产租赁住房REIT & 2026年首期每10份CNY0.5154。 & 当前市价/NAV、运营、债务与跨区域租赁住房相关性缺失。 \\
\texttt{508080} & 中金亦庄产业园REIT & 2026年首期每10份CNY0.9000。 & 当前价格、园区出租率、租户、项目现金流和流动性缺失。 \\
\texttt{508096} & 中航京能国际能源REIT & 2025年度第二次每10份CNY3.3178。 & 发电量、电价、项目现金流、债务、价格/NAV和流动性缺失。 \\
\bottomrule
\end{tabularx}
\end{exhibitbox}

\section{权益及基金候选：ETF与直接股票采用同一决策纪律}

\texttt{510880}的历史分红只能说明ETF可能分配现金，不等同于未来到账：2026年公告每10份CNY1.4300、基准日净值CNY3.1813；2025年报净值增长-0.45\%、五年45.43\%，并披露每10份分配CNY1.4200。\texttt{512890}则在2024年净值增长22.64\%的同时未作利润分配，说明合同条件、基金净值和指数股息率均不能换算为持有人下一年现金。\sourcenote{\texttt{510880} 2026年分红公告及2025年报（R10--R11）；\texttt{512890} 2024年报（R12）。}

因此不应只买ETF或只买股票。直接股票与ETF/REIT只能在同一份跨资产比较包中，按成本税费后的总回报、估值、分配与现金流质量、风险、重叠、容量和证据质量竞争；资产类别、历史分派率与预设权重均不构成胜出理由。产品的分配机制也不规定策略的现金用途：现金到账后是否保留、取出或另行考虑配置，由投资者独立决定，不改变既有总回报的计量。
